\documentclass[11pt]{article}

\usepackage[T1]{fontenc}
\usepackage[utf8]{inputenc}
\usepackage{lmodern}
\usepackage[a4paper,margin=28mm]{geometry}
\usepackage{microtype}
\usepackage{amsmath,amssymb,amsthm,mathtools}
\usepackage{enumitem}
\usepackage{booktabs,tabularx}
\usepackage[colorlinks=true,linkcolor=blue,citecolor=blue,urlcolor=blue]{hyperref}

\newtheorem{theorem}{Theorem}[section]
\newtheorem{proposition}[theorem]{Proposition}
\newtheorem{corollary}[theorem]{Corollary}
\theoremstyle{definition}
\newtheorem{definition}[theorem]{Definition}
\newtheorem{example}[theorem]{Example}
\theoremstyle{remark}
\newtheorem{remark}[theorem]{Remark}

\newcommand{\Hspace}{\mathcal H}
\newcommand{\Badm}{B_{\mathrm{adm}}}
\newcommand{\Kadm}{K_{\mathrm{adm}}}
\newcommand{\Ran}{\operatorname{Ran}}
\newcommand{\spec}{\operatorname{spec}}
\newcommand{\tr}{\operatorname{tr}}
\newcommand{\Id}{I}

\title{Finite Coherent Filters in Modal Triplet Theory:\\
\large Operator Types, Exact Kernels, and Sectoral Boundaries}
\author{Peter Nero}
\date{July 2026\\Version 3}

\begin{document}
\maketitle

\begin{abstract}
Finite coherent filtering is a useful common language for several Modal
Triplet Theory constructions, but it does not by itself derive every physical
sector in which a filter can be written. This paper gives the operator
statement a precise domain. It distinguishes an orthogonal projector, a
smooth spectral filter, and a positive measurement effect; proves positivity
and contraction of the sandwiched filter without assuming commutativity;
states the additional hypothesis needed for a modal eigenvalue formula; and
derives the Euclidean heat kernel and its sharp delta limit. The current
finite q79 Hessian supplies an exact six-mode illustration with a rank-two
coherent kernel and rank-four complement. A separate no-go theorem proves
that filtering only internal modes cannot suppress arbitrarily high
four-dimensional momenta. Gauge, gravitational, Lorentzian, and measurement
applications therefore require their own compatible operators, constraints,
and source maps. The result is a rigorous operator toolkit and a clear list
of transfer obligations, not a universal derivation of quantum field theory
or quantum gravity.
\end{abstract}

\section*{Version 3 Revision Note}
\begin{description}[leftmargin=3.2cm,style=nextline]
\item[Supersedes] Version 2, \emph{Finite Coherent Projection in Modal
Triplet Theory: A Common Architecture for Duality, Gauge, Gravity,
Quantization, Entanglement, and Measurement}.
\item[Reason] The earlier paper did not keep a sharp projector, a smooth
filter, and a positive measurement effect consistently separate. It also
risked transferring an internal spectral cutoff into four-dimensional
ultraviolet claims without a bridge theorem.
\item[Resolution] The operator variables and spectra are now typed
explicitly. The central analytic theorem is corrected, the q79 finite
example is imported at its established tier, and an internal-to-spacetime
no-go theorem is added.
\item[Retained result] The positive-contraction filter, its local and
spectral representations, and the heat-kernel delta limit remain valid under
the stated assumptions.
\item[Remaining boundary] Physical gauge, gravitational, Lorentzian, and
measurement uses still require selected same-source operators and
compatibility theorems. No four-dimensional UV completion follows from the
internal projector alone.
\end{description}

\tableofcontents

\section{Why the distinctions matter}

The phrase ``finite coherent projection'' originally served as an umbrella
for three operations. They are related, but they answer different questions.

\begin{enumerate}
\item A \emph{sharp projector} decides membership in a closed subspace.
\item A \emph{smooth filter} attenuates spectral components by different
amounts.
\item A \emph{measurement effect} assigns an outcome probability in a
specified quantum model.
\end{enumerate}

Conflating them creates immediate errors. A smooth heat filter is usually not
idempotent, so it is not a projector. A positive filter may be used as one
effect of a measurement, but it does not specify the state update or the
other outcomes. An internal filter can remove internal modes while leaving
four-dimensional momentum untouched.

The corrected role of this paper is consequently modest but useful. It
provides the common operator calculus, gives an exact finite MTT
instantiation, and lists the extra data needed before the same notation can
be used in a physical sector.

\section{Three operator types}

\begin{definition}[Sharp projector]
An orthogonal projector on a Hilbert space \(\Hspace\) is a bounded operator
\[
 P=P^\ast=P^2.
\]
It represents the exact closed subspace \(\Ran P\).
\end{definition}

\begin{definition}[Smooth spectral filter]
Let \(A\) be a nonnegative self-adjoint operator on \(\Hspace\). For a bounded
Borel function \(f:[0,\infty)\to[0,1]\), functional calculus defines
\[
 F=f(A), \qquad 0\leq F\leq \Id.
\]
Unless \(f\) takes only the values zero and one on \(\spec A\), \(F\) is not
a projector.
\end{definition}

\begin{definition}[Effect and instrument]
An effect is an operator \(E\) satisfying \(0\leq E\leq\Id\). A finite POVM
is a family \(\{E_i\}\) with \(\sum_iE_i=\Id\). An instrument additionally
specifies completely positive maps whose outcome probabilities and
post-measurement states agree with the effects.
\end{definition}

Thus every projector is an effect, but not every effect is a projector.
Neither object alone specifies a measurement instrument. This distinction is
important in MTT because admissibility, attenuation, and record formation
belong to different logical layers.

\section{Specified filter data}

The abstract data used in the rest of the paper are:
\[
(\Hspace,A,P,\chi,\tau),
\]
where
\begin{itemize}
\item \(\Hspace\) is a complex Hilbert space;
\item \(A\geq0\) is self-adjoint, with its domain fixed;
\item \(P\) is an orthogonal projector;
\item \(\chi:[0,\infty)\to[0,1]\) is bounded and Borel measurable;
\item \(\tau\geq0\) is a declared filter parameter.
\end{itemize}

Define
\begin{equation}
 G_\tau=\chi(A)e^{-\tau A/2},
 \qquad
 \Badm=P\,G_\tau^2P
       =P\chi(A)e^{-\tau A}\chi(A)P.
 \label{eq:Badm}
\end{equation}
The symbol \(A\) is not universal. On a compact internal manifold \(Y\) it
may be a nonnegative elliptic operator on \(L^2(Y,E)\), in which case its
spectrum is discrete under standard hypotheses. On Euclidean spacetime it
may be \(-\Delta\). On a Cauchy slice it may be a positive spatial operator.
These are different choices on different Hilbert spaces and may not be
silently identified.

\section{The finite-filter theorem}

\begin{theorem}[Sandwiched finite filter]
\label{thm:filter}
For the data above, no commutation assumption between \(P\) and \(A\) is
needed to conclude
\[
 0\leq \Badm\leq\Id,
 \qquad
 \Ran\Badm\subseteq\Ran P.
\]
If in addition \([P,A]=0\), and
\[
 A\phi_n=\lambda_n\phi_n,\qquad
 P\phi_n=p_n\phi_n,\qquad p_n\in\{0,1\},
\]
then
\[
 \Badm\phi_n
 =p_n\chi(\lambda_n)^2e^{-\tau\lambda_n}\phi_n.
\]
\end{theorem}

\begin{proof}
Because \(\chi(A)\) and \(e^{-\tau A/2}\) are commuting bounded functions of
\(A\), \(G_\tau\) is a positive contraction. Hence
\[
 \Badm=(G_\tau P)^\ast(G_\tau P)\geq0
\]
and
\[
 \|\Badm\|\leq\|G_\tau P\|^2\leq1.
\]
The outer factor \(P\) places the range in \(\Ran P\). Under the additional
commutation hypothesis, \(A\) and \(P\) admit the displayed common modal
action, and functional calculus gives the eigenvalue formula.
\end{proof}

\begin{corollary}[When the filter is actually a projector]
\label{cor:idempotence}
In the commuting pure-point setting of \autoref{thm:filter}, \(\Badm\) is a
projector exactly when every retained weight
\[
 p_n\chi(\lambda_n)^2e^{-\tau\lambda_n}
\]
is zero or one. In particular, for \(\tau>0\), every retained mode with
\(\lambda_n>0\) has weight strictly below one whenever
\(\chi(\lambda_n)\neq0\).
\end{corollary}

This is why \(\Badm\) should be called a finite coherent filter, not the
coherent projector. The sharp object is \(P\); the attenuation is carried by
\(G_\tau\).

\subsection{Kernel representation}

When \(\Badm\) has an integral kernel,
\[
 (\Badm\psi)(x)=\int \Kadm(x,y)\psi(y)\,dy.
\]
In a commuting discrete eigenbasis,
\begin{equation}
 \Kadm(x,y)=
 \sum_n p_n\chi(\lambda_n)^2e^{-\tau\lambda_n}
 \phi_n(x)\overline{\phi_n(y)}.
 \label{eq:spectral-kernel}
\end{equation}
The coordinate-space kernel and spectral sum are two representations of the
same operator. Calling one a ``particle shadow'' and the other a ``wave
shadow'' is an MTT interpretation of that dual representation. The equation
alone does not derive complementarity, the Born rule, or detector dynamics.

\section{The exact heat-kernel benchmark}

\begin{theorem}[Euclidean heat kernel and delta limit]
\label{thm:heat}
Let \(\Hspace=L^2(\mathbb R^d)\), \(A=-\Delta\), \(P=\Id\), and
\(\chi=1\). For \(\tau>0\),
\[
 \Badm=e^{\tau\Delta}
\]
has kernel
\begin{equation}
 K_\tau(x,y)=
 (4\pi\tau)^{-d/2}
 \exp\!\left(-\frac{|x-y|^2}{4\tau}\right).
 \label{eq:heat-kernel}
\end{equation}
Moreover \(e^{\tau\Delta}\psi\to\psi\) strongly in \(L^2\), and
\(K_\tau(x,\cdot)\to\delta_x\) in the distributional sense as
\(\tau\downarrow0\).
\end{theorem}

\begin{proof}
The Fourier transform diagonalizes \(-\Delta\), so
\[
 \widehat{e^{\tau\Delta}\psi}(k)=e^{-\tau|k|^2}\widehat\psi(k).
\]
Fourier inversion gives \eqref{eq:heat-kernel}. Dominated convergence proves
strong \(L^2\) convergence. Pairing the normalized Gaussian with a test
function and rescaling \(y=x+\sqrt{\tau}\,z\) proves the distributional
limit.
\end{proof}

This theorem is exact, but it is a Euclidean spacetime benchmark. It is not
obtained merely by selecting internal q79 modes. It also says nothing by
itself about Lorentzian support, reflection positivity, unitarity, gauge
identities, or removal of a quantum-field-theory regulator.

\section{The exact finite q79 illustration}

The current MTT Foundation establishes, on its selected finite q79
flat-symbol carrier, a normalized Hessian
\[
 A_{\mathrm{fin}}=\kappa_{\mathrm{fin}}^{-1}H_{\mathrm{fin}}
\]
with
\[
 \spec A_{\mathrm{fin}}
 =\{0\text{ with multiplicity }2,\;
    1\text{ with multiplicity }4\}.
\]
It also identifies the rank-two Haar projector \(P_{\mathrm{Haar}}\) with
the zero eigenspace and proves compatibility with the shared finite line and
the \(1+2+3\) sector projectors. Those are imported Foundation results; this
paper does not claim ownership of them.

Writing \(Q=\Id-P_{\mathrm{Haar}}\), functional calculus is completely
explicit:
\begin{align}
 e^{-\tau A_{\mathrm{fin}}}
 &=P_{\mathrm{Haar}}+e^{-\tau}Q,\\
 \chi(A_{\mathrm{fin}})
 &=\chi(0)P_{\mathrm{Haar}}+\chi(1)Q.
\end{align}
Consequently, with \(P=\Id\),
\begin{equation}
 \Badm=
 \chi(0)^2P_{\mathrm{Haar}}
 +\chi(1)^2e^{-\tau}Q.
 \label{eq:q79-filter}
\end{equation}
If instead \(P=P_{\mathrm{Haar}}\), then
\[
 \Badm=\chi(0)^2P_{\mathrm{Haar}}.
\]

Equation \eqref{eq:q79-filter} is a genuine exact finite filter. It
distinguishes a two-dimensional coherent kernel from a four-dimensional
penalized complement. It does not select a physical value of \(\tau\), a
continuum HYM operator, a four-dimensional propagator, or a detector effect.
Those identifications require further source maps.

\section{Why internal filtering is not a spacetime UV cutoff}

\begin{theorem}[Internal-filter UV no-go]
\label{thm:no-go}
Let \(\Hspace=\Hspace_4\otimes\Hspace_{\mathrm{int}}\). Let \(A_4\) be an
unbounded nonnegative self-adjoint operator on \(\Hspace_4\) whose spectral
subspace above every finite threshold is nonzero. If
\(B_{\mathrm{int}}\neq0\) is bounded, then
\[
 R=\Id_4\otimes B_{\mathrm{int}}
\]
does not annihilate or attenuate all sufficiently high \(A_4\)-modes.
Therefore an internal finite-rank projector or smooth filter, acting alone,
is not a four-dimensional ultraviolet regulator.
\end{theorem}

\begin{proof}
For any cutoff \(\Lambda\), choose a nonzero
\(u_\Lambda\) in the spectral subspace
\(\mathbf 1_{(\Lambda,\infty)}(A_4)\Hspace_4\). Since
\(B_{\mathrm{int}}\neq0\), choose \(v\) with
\(B_{\mathrm{int}}v\neq0\). Then
\[
 R(u_\Lambda\otimes v)
 =u_\Lambda\otimes B_{\mathrm{int}}v\neq0.
\]
The four-dimensional spectral label of \(u_\Lambda\) is unchanged, and
\(\Lambda\) was arbitrary.
\end{proof}

The theorem leaves open a useful route. A derived four-dimensional filter
could arise after reduction if a proved intertwiner maps internal geometry
to a spacetime operator \(A_4\), including its normalization and constraints.
But that map is precisely an additional theorem. It cannot be replaced by
using the same letter \(A\) in both settings.

\section{What sectoral promotion requires}

\subsection{Gauge theory}

A gauge-sector filter must be defined on a gauge-covariant complex or on the
physical quotient. At minimum one needs a relation such as
\[
 [F,C_a]=0
\]
on a common invariant domain for the constraints \(C_a\), or a chain-map
identity with the BRST differential. A generic spectral filter can violate
Ward or Slavnov--Taylor identities. The presence of a Lens-like quotient in
the MTT vocabulary does not establish these analytic identities.

\subsection{Gravity}

For linearized gravity one may specify a transverse-traceless projector and
a positive spatial or Euclidean operator. Nonlinear gravity additionally
requires diffeomorphism constraints, lapse and shift treatment, boundary
terms, and a selected action. A filter on internal bundle modes does not
automatically define a diffeomorphism-compatible graviton propagator.

\subsection{Lorentzian evolution}

The heat kernel in \autoref{thm:heat} is Euclidean. Replacing \(-\Delta\) by a
hyperbolic wave operator does not automatically produce a positive,
retarded, or causal filter. A bounded spatial filter can be used on Cauchy
data, but if it is nonlocal it may still alter support properties. A physical
claim therefore needs a retarded kernel, a support estimate, or a
preparation/readout interpretation.

\subsection{Measurement}

By \autoref{thm:filter}, \(\Badm\) is an effect. It can define the binary POVM
\(\{\Badm,\Id-\Badm\}\) in standard quantum mechanics. This does not select
an instrument, derive quadratic probabilities from pre-quantum MTT data, or
pick one ontic outcome. Those are separate source and dynamics questions.

\subsection{Entanglement}

On a bipartite Hilbert space, a joint filter need not factor:
\[
 B_{AB}\neq B_A\otimes B_B.
\]
That observation permits correlated or entangled coherent sectors, but
nonfactorization alone does not derive a particular state, Bell
correlations, or no-signaling. Those properties must be checked in the
chosen state and observable algebra.

\section{Circle, Lens, and Nil after reconciliation}

The safest use of Circle--Lens--Nil in this paper is an operator and
filtration language:
\begin{description}[leftmargin=2.4cm,style=nextline]
\item[Circle] phase or holonomy data carried by a specified line bundle and
connection;
\item[Lens] quotient, identification, or invariant-sector data implemented
by a specified projector or group action;
\item[Nil] survivor, shear, or non-semisimple data implemented by a
specified operator and domain.
\end{description}

This is not a claim that the physical compactification is the literal nested
manifold ``circle inside Lens inside Nil.'' The selected q79 Fu--Yau branch
and the finite shared line have their own topology. Equality of suggestive
roles does not establish equality of bundles, connections, Hessians, or
spectra.

\section{Status ledger}

\begin{center}
\begin{tabularx}{\textwidth}{@{}
>{\raggedright\arraybackslash}p{0.25\textwidth}
>{\raggedright\arraybackslash}p{0.19\textwidth}
>{\raggedright\arraybackslash}X@{}}
\toprule
Statement & Status & Evidence or missing object\\
\midrule
Positive-contraction theorem & Exact &
\autoref{thm:filter}, by functional calculus.\\
Projector/filter/effect distinction & Exact &
Definitions and \autoref{cor:idempotence}.\\
Euclidean Gaussian kernel and delta limit & Exact &
\autoref{thm:heat}.\\
Finite q79 two-plus-four filter & Exact, imported &
Foundation finite Hessian and Haar-projector theorem.\\
Internal filter is not a 4D UV cutoff & Exact &
\autoref{thm:no-go}.\\
Physical continuum q79 filter & Open &
Selected nonflat HYM operator, domain, spectrum, and normalization.\\
Gauge-compatible filter & Conditional &
Constraint or BRST chain-map theorem is required.\\
Gravity-compatible filter & Conditional &
Selected action, constraints, causal kernel, and continuum limit are
required.\\
Measurement model & Conditional/open &
POVM use is standard; selected MTT instrument and general Born source remain
separate.\\
Universal projection-shadow theorem & Withdrawn &
One operator template does not derive all sectoral physics.\\
\bottomrule
\end{tabularx}
\end{center}

\section{Discussion}

The corrected result is stronger than a broad analogy because it is
falsifiable at each transfer. Once \(\Hspace\), \(A\), \(P\), \(\chi\), and
\(\tau\) are declared, positivity, idempotence, spectral weights, and kernel
limits can be checked. When a proposed application changes Hilbert space,
constraints, signature, or observable algebra, the required bridge becomes
visible rather than being hidden in notation.

The q79 finite carrier is particularly informative. It shows that MTT
currently has an exact finite spectral decomposition, not merely a verbal
picture of coherence. At the same time, \autoref{thm:no-go} identifies the
precise limit of that success: the finite internal spectrum is not yet a
four-dimensional UV theorem. A future derivation must transport the selected
geometry into a physical spacetime operator while preserving the relevant
connections, constraints, Hessians, and causal structure.

This separation also clarifies the role of finite cutoff exactness. If the
selected physical object is itself a finite projected algebra, its trace and
functional calculus are exact at that tier. Exactness of the finite object
does not imply that it approximates a continuum theory with zero error. A
continuum claim needs either an equivalence theorem or explicit convergence
and error bounds.

\section{Conclusion}

Finite coherent filtering survives corpus reconciliation as a rigorous and
useful MTT tool. The central operator is a positive contraction; it has
well-defined modal and kernel representations; its Euclidean heat-kernel
specialization converges to a delta distribution; and the current finite q79
Hessian gives an exact two-plus-four example.

The same analysis closes an important loophole. An internal projector does
not suppress high four-dimensional momentum merely because both operations
are called filters. Physical promotion requires a selected operator on the
correct Hilbert space together with gauge, gravitational, Lorentzian, or
measurement compatibility as appropriate. That is the right frontier for
this architecture.

\begin{thebibliography}{99}

\bibitem{MTTFoundation}
P.~Nero,
\emph{Modal Triplet Theory: Foundations},
current revised MTT paper corpus, 2026.

\bibitem{MTTQGOverview}
P.~Nero,
\emph{A Projection-First Reframing of Quantum Gravity},
current revised MTT paper corpus, 2026.

\bibitem{MTTSPT}
P.~Nero,
\emph{A Conditional Euclidean Transverse-Traceless Filter Model in Modal
Triplet Theory},
current revised MTT paper corpus, 2026.

\bibitem{ReedSimon}
M.~Reed and B.~Simon,
\emph{Methods of Modern Mathematical Physics I: Functional Analysis},
Academic Press, 1980.

\bibitem{Davies}
E.~B.~Davies,
\emph{Heat Kernels and Spectral Theory},
Cambridge University Press, 1989.

\bibitem{Holevo}
A.~S.~Holevo,
\emph{Probabilistic and Statistical Aspects of Quantum Theory},
North-Holland, 1982.

\bibitem{Henneaux}
M.~Henneaux and C.~Teitelboim,
\emph{Quantization of Gauge Systems},
Princeton University Press, 1992.

\bibitem{BarGinouxPfaeffle}
C.~B\"ar, N.~Ginoux, and F.~Pf\"affle,
\emph{Wave Equations on Lorentzian Manifolds and Quantization},
European Mathematical Society, 2007.

\end{thebibliography}

\end{document}
